\section*{Ex.~2.7 (22.5 pts) --- Three papers from the ``Scheduling'' group}

\subsection*{Paper 2-3 --- Waldspurger and Weihl, \emph{Lottery Scheduling} (OSDI 1994)}

\paragraph{Summary.}
Lottery scheduling represents resource rights as \emph{tickets} and settles each quantum by
drawing one at random: a client holding $t$ of $T$ tickets wins with probability $p = t/T$, so
throughput tracks the allocation, no non-zero holding can starve, and a reallocation takes effect
at the very next decision. The abstractions built on the draw matter more than the draw itself.
\emph{Ticket transfers} hand a blocked client's rights to whoever is computing on its
behalf, dissolving priority inversion much as priority inheritance does. \emph{Currencies}
denominate tickets locally and convert through an exchange rate into a conserved base currency,
so inflation inside one trust boundary cannot leak out. \emph{Compensation tickets} multiply a
client's funding by $1/f$ when it used only a fraction $f$ of its quantum, so a thread that
yields early for I/O is not quietly punished. A Mach 3.0 prototype (25\,MHz MIPS DECStation,
100\,ms quantum) reports a $2\!:\!1$ Dhrystone allocation yielding $25378$ against $12619$
iterations/s ($2.01\!:\!1$), an $8\!:\!3\!:\!1$ database allocation whose response times while the
best-funded client is active ($17.19$, $43.19$ and $132.20$\,s) give relative speeds of
$7.69\!:\!2.51\!:\!1$, and currencies insulating one user's
inflation from another (aggregate ratio $1.01\!:\!1$ before, $1.00\!:\!1$ after).

\paragraph{Critical judgment.}
The price of randomisation is variance, and the paper meets that objection rather than ducking
it: the one bad run, a $10\!:\!1$ allocation observed at $13.42\!:\!1$, is named as the exception
to ratios otherwise ``close to their corresponding allocations'', and Figure~5 goes on to plot
8-second windows inside a 200-second run. The weak step is the last one, where fairness at the
timescale the introduction cares about --- interactive control ``at a time scale of milliseconds
to seconds'' --- is obtained by argument instead of measurement: ``if a scheduling quantum of 10
milliseconds were used instead of the 100 millisecond Mach quantum, the same degree of fairness
would be observed over a series of subsecond time windows''. Nothing was ever run at 10\,ms. The
warrant offered for a fine quantum being cheap --- ``a tree-based lottery need only generate a
random number and perform $\lg n$ additions and comparisons'' --- describes an implementation
that was never built, whereas every measurement in the paper comes from the unoptimised
list-based prototype.

The overhead claim is thinner than its conclusion. The prototype ran $2.7\%$ slower on three
Dhrystone tasks and $0.8\%$ slower on eight, and the authors honestly note the run-to-run
standard deviation was as large as the effect; the one significant measurement, $1.7\%$
\emph{faster} on the database server, is attributed in a footnote to incidental cache and TLB
locality. ``Efficient lottery scheduling does not pose any challenging problems'' therefore rests
on data that mostly cannot separate the two schedulers. Currencies, finally, are what make
inflation safe, and \S5.5 does demonstrate the barrier containing one user's inflation; what is
missing is the control over \emph{who} may inflate a currency, since the access-control lists
that would supply it are named only as something ``a complete lottery scheduling system should''
provide.

\subsection*{Paper 2-4 --- Waldspurger and Weihl, \emph{Stride Scheduling} (MIT/LCS/TM-528, 1995)}

\paragraph{Summary.}
Stride scheduling is the same authors' deterministic replacement for their own dice. Each client
holds a \emph{stride} inversely proportional to its tickets and a \emph{pass} value; the
scheduler always runs the minimum-pass client and advances that client's pass by its stride, in
$O(\lg n_c)$ time with a priority queue. A partial quantum advances the pass by only
$f \cdot \mathit{stride}$, while joins, leaves and ticket changes are handled by saving and
rescaling the client's unconsumed remainder against a smoothly advancing global pass. Because
each client's state is independent, the transfers, inflation and currencies of paper 2-3 all
survive on the deterministic substrate. The error
bounds are the argument: lottery scheduling's absolute error grows as $O(\sqrt{n_a})$ in the
number of allocations, whereas stride's pairwise relative error never exceeds one quantum
whatever $n_a$ is. Absolute error can still reach $O(n_c)$ under skew, $101$ clients at
$100\!:\!1\!:\!\dots\!:\!1$ giving an error of $50$, which motivates \emph{hierarchical} stride
scheduling: a balanced tree of aggregated groups bounding absolute error to $O(\lg n_c)$, which
measured $4.5$ on that example. In simulation at $19\!:\!1$ over one million allocations, the
one-ticket client saw response times of $\mu = 20.00$, $\sigma = 0.01$ quanta under stride
against $\mu = 20.13$, $\sigma = 19.64$ and a range of $1$--$194$ quanta under lottery. Linux
1.1.50 prototypes held CPU ratios within $1\%$ of ideal ($3\!:\!1$ gave $3.001\!:\!1$) and UDP
ratios within $5\%$, in under 300 lines each.

\paragraph{Critical judgment.}
The head-to-head case against lottery scheduling is made entirely in simulation. The prototypes
are measured against the \emph{ideal} ratio and against stock Linux, never against a
lottery-scheduled kernel, so the headline claim of significantly improved accuracy over lottery
is nowhere demonstrated on hardware. They also leave the modularity argument unexercised:
Section 3 shows how transfers and currencies would be built, but the CPU prototype implements
neither, so the abstractions that justify the design are never run. The overhead figure is softer
than it reads, because the authors concede their prototype does a linear scan rather than the
$O(\lg n_c)$ structure the complexity analysis assumes; the cheap measurement and the good bound
describe different programs.

Against that, the paper is unusually candid and criticism should not pretend otherwise. It
reports that the network result required lowering a hard-coded delay constant in \texttt{ttcp},
without which ratios were ``consistently lower than specified''; it volunteers that its core
algorithm turned out to be nearly VirtualClock and WFQ rediscovered; and its conclusion concedes
that lottery scheduling is ``conceptually simpler'' and ``effectively stateless'' while stride
demands careful state fixups on every dynamic change. History settled the trade: Linux's CFS and
its EEVDF successor both schedule by picking the minimum-virtual-time task and advancing it by a
weight-scaled increment, stride's rule under another name, while no production kernel ever
shipped the lottery.

\subsection*{Paper 2-1 --- Buttazzo, \emph{Rate Monotonic vs.\ EDF: Judgment Day} (Real-Time Systems 29(1), 2005)}

\paragraph{Summary.}
This is a rebuttal rather than a new result. Buttazzo takes the five claims habitually used to
prefer Rate Monotonic to Earliest Deadline First (simpler implementation, lower overhead, easier
analysis, more predictable overload, less jitter) and answers each with theory, counterexamples
and simulation. He grants RM one real point: on a commercial kernel with at most $256$ priority
levels, a new deadline falling between two adjacent levels can force every active deadline to be
remapped, and only RM admits $O(1)$ insertion via one FIFO per level. Written natively, the two
differ only in the ready-queue sort key. On overhead he argues
the accounting is backwards: EDF pays one deadline computation per job, but RM forces far more
preemptions, and his simulations show preemption counts rising roughly linearly with load under
RM while \emph{falling} under EDF at high load, because a long job holding the earliest deadline
can push a successor past an arrival and avoid a preemption outright. On analysis, EDF is exact
in $O(n)$ at $U \le 1$ when deadlines equal periods, while RM's $\ln 2 \approx 0.69$ bound needs
pseudo-polynomial response-time analysis to become exact.
Under permanent overload EDF rescales every period to $T_i U$ while RM can
block low-priority tasks outright; under transient overload a $1.5$-unit overrun in the
\emph{highest}-rate task makes $\tau_2$ miss rather than the longest-period $\tau_4$, refuting
the belief that RM fails predictably from the bottom up. A three-task example gives jitters of
$0, 2, 8$ under RM against $1, 2, 3$ under EDF.

\paragraph{Critical judgment.}
The method is asymmetric, and the paper is more careful about it than its title. A counterexample
can refute ``RM always wins'' but cannot establish ``EDF wins'', and Buttazzo keeps to the weaker
claim: the jitter example ``does not prove that EDF always introduces less jitter\ldots but just
confutes the common belief'', and the conclusion says only that RM's advertised properties ``do
not hold in general''. What overshoots is the register, not the claims: a Judgment Day that
delivers a rebuttal. The overhead case is the weakest link: Section~3 counts preemptions, never measures a
context switch on real hardware, and nets the saved switches against EDF's per-job deadline
update by assertion alone (``EDF introduces less runtime overhead than RM, when context switches
are taken into account''), so the verdict rests on a proxy counted over synthetic task sets. Two
things that would decide the industrial argument are absent: multiprocessors, and the
certification ecosystem that keeps fixed priorities in use.

He does not, however, stack the deck. He grants RM the $O(1)$ ready queue and the cheaper exact
test when deadlines are shorter than periods, and grants that RM confines an overrun to tasks of
lower priority while EDF confines it nowhere --- adding at once that the property ``can be of
little use if we do not know a priori which task is going to overrun''. He concludes that under
\emph{permanent} overload deciding which behaviour is better is ``highly application dependent'',
and that both algorithms are ``not very well suited'' to overload and jitter control at all, the
remedy being resource reservation. That concession aged best:
Linux's
\texttt{SCHED\_DEADLINE} (merged in 3.14, 2014) ships EDF together with the Constant Bandwidth
Server, the Abeni--Buttazzo mechanism of his own Section 7.3, rather than EDF alone.

\medskip
Two of the three are acts of replacement: 2-4 unseats its own predecessor's randomness and 2-1
unseats static rate priorities, both on the ground that a scheduler should carry the extra state
needed to decide accurately --- exactly the state 2-3 had traded away for a stateless draw.

\subsection*{References}

\begin{enumerate}
  \item[{[1]}] C.~A.~Waldspurger and W.~E.~Weihl. ``Lottery Scheduling: Flexible
    Proportional-Share Resource Management.'' \emph{Proceedings of the First USENIX Symposium on
    Operating Systems Design and Implementation (OSDI)}, Monterey, CA, November 1994.
    (Paper list ID 2-3.)
  \item[{[2]}] C.~A.~Waldspurger and W.~E.~Weihl. ``Stride Scheduling: Deterministic
    Proportional-Share Resource Management.'' Technical Memorandum MIT/LCS/TM-528, MIT Laboratory
    for Computer Science, June 1995. (Paper list ID 2-4.)
  \item[{[3]}] G.~C.~Buttazzo. ``Rate Monotonic vs.\ EDF: Judgment Day.'' \emph{Real-Time
    Systems}, 29(1):5--26, 2005. (Paper list ID 2-1.)
\end{enumerate}
